% SPDX-License-Identifier: MIT
\chapter{Layouts}
\label{ch:layouts}

There are five layout modes, selected with keys \keys{1}--\keys{5}, from the
Style panel, or with \opt{--mode} on the command line. Independently of the mode,
the \emph{branch mode} decides how the along-axis coordinate is derived.

\section{Branch modes}

\begin{center}
\begin{tabular}{@{}llp{0.5\linewidth}@{}}
\toprule
\textbf{Branch mode} & \textbf{Key} & \textbf{Along-axis coordinate} \\
\midrule
Phylogram              & \keys{Ctrl+1} & Cumulative branch length. Distances are meaningful. \\
Cladogram, aligned tips& \keys{Ctrl+2} & Depth, with every tip on the outer edge. Topology only. \\
Cladogram, by level    & \keys{Ctrl+3} & Node depth in levels. \\
\bottomrule
\end{tabular}
\end{center}

A phylogram shows evolutionary distance; a cladogram shows only branching order.
Use a cladogram when branch lengths are absent or meaningless, and say which you
used in the figure legend --- readers cannot tell by looking.

\section{Rectangular}

The default, and the right choice for most figures. Branch length runs along $x$,
tips are stacked down $y$.

\galleryfig{layout-rectangular-phylogram}{Rectangular phylogram. The scale bar at
the lower left gives the units of the horizontal axis.}{1.0}

\galleryfig{layout-rectangular-cladogram}{Rectangular cladogram with aligned
tips. Horizontal position now carries no distance information.}{1.0}

\section{Slanted}

Branches are drawn as straight diagonals rather than as right angles. Compact and
readable for small trees; harder to follow for large ones, because it is less
obvious which ancestor a branch descends from.

\galleryfig{layout-slanted}{Slanted layout.}{1.0}

\section{Circular}

Tips are placed around an arc. This is the layout that makes a large tree fit on
a page, and it is the natural home for multiple annotation rings.

Three options shape it:

\begin{center}
\begin{tabular}{@{}lll@{}}
\toprule
\textbf{Option} & \textbf{Default} & \textbf{Effect} \\
\midrule
\opt{arc}          & 350\textdegree & Angular span. 360\textdegree{} closes the circle. \\
\opt{inner\_radius}& 0.1            & Hole at the centre, as a fraction of the radius. \\
\opt{start\_angle} & $-90$\textdegree & Where the first tip sits. \\
\bottomrule
\end{tabular}
\end{center}

Leaving a small gap --- the default 350\textdegree{} --- makes it clear where the
tip order begins and ends. Closing to a full circle looks tidier but loses that
cue.

The gap does a second job that is easy to miss: it is the only place a track
caption can go. A rectangular tree captions its columns in a header row above
them and grows the page to fit; a fan has no room above its rings, so the
captions are written into the wedge the arc leaves open. Ask for several rings
with long names and the wedge fills up, at which point MakeYourTree reports the
captions it could not place rather than drawing them across the data. Widening
the gap --- \opt{arc} at 340\textdegree{} or 320\textdegree{} --- is the fix;
the legend names every track either way.

\galleryfig{layout-circular}{Circular fan with the default 350\textdegree{} arc.
Labels are rotated tangentially and flipped on the left-hand side so that none of
them is upside down.}{0.75}

\galleryfig{layout-circular-full}{A full 360\textdegree{} circle with a larger
inner radius.}{0.75}

\begin{quote}
\textbf{A gap you should not try to fix.} In a circular \emph{phylogram}, tips sit
at different radii because they are at different evolutionary distances from the
root. A ragged inner edge against an annotation ring is therefore correct, not a
bug. Turn on \textbf{Align tips} (\keys{A}) if you want the ring flush.
\end{quote}

\section{Radial}

Angle comes from tip order and radius from branch length, but the tree is drawn
outward from a central root without the concentric-arc branch style of the
circular mode.

\galleryfig{layout-radial}{Radial layout.}{0.75}

\section{Unrooted}

An unrooted layout makes no claim about which node is ancestral. Two algorithms
are available.

\begin{description}[style=nextline,leftmargin=0pt,labelindent=0pt,labelwidth=0pt]
  \item[Equal-angle] Felsenstein's original method. Each subtree receives a wedge
    proportional to the number of tips it contains. Fast and always planar.
  \item[Equal-daylight] Iteratively adjusts the wedges so the blank space between
    subtrees is evened out. Slower, usually easier to read.
\end{description}

Branch lengths are preserved exactly in both, and both are planar by
construction.

\galleryfig{layout-unrooted-equal-angle}{Unrooted, equal-angle.}{0.85}

\galleryfig{layout-unrooted-equal-daylight}{Unrooted, equal-daylight. The
subtrees are spread to even out the gaps between them.}{0.85}

\subsection{Annotation tracks in unrooted layouts}
\label{sec:unrooted-tracks}

\begin{quote}
\textbf{Known limitation.} An unrooted layout has no tip ordering along a line or
a circle, so band space has no faithful meaning in it. Tracks are currently drawn
as a column beside the figure, in tip order, \emph{aligned to nothing}. The
composition step emits the diagnostic
\cmd{compose.approximate-band-space} to record this, but the interface does not
yet present it prominently.

Do not use annotation tracks with an unrooted layout in a published figure
without saying so. Use a rooted layout instead.
\end{quote}

\section{Spacing and labels}

\begin{center}
\begin{tabular}{@{}lp{0.6\linewidth}@{}}
\toprule
\textbf{Setting} & \textbf{Effect} \\
\midrule
\opt{row\_spacing}   & Distance between adjacent tip rows. The main control over figure height. \\
\opt{width}, \opt{height} & Page size in scene units. \\
\opt{align\_tips}    & Flush all tip labels to a common edge. \\
\opt{guide\_lines}   & Dotted leaders from each tip to its label. Only meaningful with aligned tips. \\
\opt{show\_tip\_labels} & Turn tip labels off entirely --- often necessary above a few hundred tips. \\
\opt{max\_label\_width} & Truncate long labels with an ellipsis. \\
\bottomrule
\end{tabular}
\end{center}

\galleryfig{feature-align-tips}{Tips aligned with guide lines. This is almost
always what you want once an annotation track is attached, because it makes the
rows line up visually with the track cells.}{0.95}
